% META: {"id": "TNSI-ARCH-CO-007", "chapitre": "TNSI-ARCHITECTURES-MATERIELLES-SY", "type_objet": "corrige", "exercice_ref": "TNSI-ARCH-EX-007", "capacites": ["T-ARCH-04B"], "capacites_codes": ["C7"], "mode_creation": "ex_nihilo", "sources_inspiration": [], "status": "generated"}
% BEGIN-VERIFY
% p, q, e = 11, 13, 7
% n, phi = p * q, (p - 1) * (q - 1)
% assert (n, phi) == (143, 120)
% d = pow(e, -1, phi)
% assert d == 103 and (e * d) % phi == 1
% assert pow(42, e, n) == 81
% assert pow(81, d, n) == 42
% assert pow(pow(5, e, n), d, n) == 5
% assert pow(pow(100, e, n), d, n) == 100
% def chiffrer_xor(message, cle):
%     return bytes(octet ^ cle for octet in message)
% assert chiffrer_xor(b"BONJOUR", 42) == b"hed`e\x7fx"
% assert chiffrer_xor(chiffrer_xor(b"BONJOUR", 42), 42) == b"BONJOUR"
% assert [k for k in range(2, n) if pow(k, e, n) == 81] == [42]
% assert n == 143 and 11 * 13 == 143
% assert phi == (11 - 1) * (13 - 1)
% assert len(str(n)) == 3
% END-VERIFY
\begin{corrige}{TNSI-ARCH-EX-007}
\textbf{1.} $n=11\times 13=143$ et $\varphi=10\times 12=120$. On cherche $d$ tel
que $7d\equiv 1\pmod{120}$ : $d=103$ convient, car $7\times 103=721=6\times
120+1$. La clé publique est donc $(143\,;\,7)$ et la clé privée $103$.

Cette relation est le c\oe{}ur du procédé : c'est elle qui garantit que déchiffrer
défait exactement ce que chiffrer a fait, autrement dit que
$(m^e)^d\equiv m \pmod n$.

\textbf{2.} L'entier transmis est $42^7 \bmod 143 = 81$. C'est ce nombre, et lui
seul, qui circule sur le réseau.

\textbf{3.} $81^{103}\bmod 143 = 42$ : le serveur retrouve la clé. La
vérification tient aussi pour $5$, qui se chiffre en $47$ et se déchiffre en $5$,
et pour $100$, qui se chiffre en $100$ — un point fixe, curiosité due à la
petitesse des nombres choisis et qui disparaît avec des clés réalistes.

\textbf{4.} Les deux machines partagent le secret $42$. On l'emploie comme clé du
chiffrement symétrique :
\begin{python}
def chiffrer_xor(message, cle):
    return bytes(octet ^ cle for octet in message)

cache = chiffrer_xor(b"BONJOUR", 42)   # b"hed`e\x7fx"
chiffrer_xor(cache, 42)                # b"BONJOUR"
\end{python}
Appliquer deux fois le \lstinline{ou exclusif} avec la même clé redonne le
message : c'est bien un chiffrement symétrique.

\textbf{5.} Non, l'espion ne peut pas. Il dispose de $81$, de $n=143$ et de
$e=7$, mais pas de $d$. Pour obtenir $d$ il lui faudrait $\varphi$, donc
connaître $p$ et $q$, donc \textbf{factoriser} $n$.

Sur $143$, c'est l'affaire de quelques secondes — on peut même, ici, essayer les
$141$ clés possibles et constater que $42$ est la seule dont le chiffré vaut
$81$. Mais avec des nombres premiers de plusieurs centaines de chiffres, aucun
algorithme connu ne factorise $n$ en un temps raisonnable, alors que le calcul
de $m^e \bmod n$ reste immédiat. Toute la sécurité tient dans cet écart entre une
opération facile et son inverse.

\textbf{6.} Deux raisons, et les deux comptent.

D'abord le \textbf{coût}. Une exponentiation modulaire sur des nombres de mille
bits est plusieurs centaines de fois plus lente qu'un chiffrement symétrique sur
la même quantité de données. Chiffrer ainsi une vidéo entière saturerait le
serveur, alors que chiffrer une seule clé de quelques dizaines d'octets ne coûte
rien.

Ensuite la \textbf{nature des données}. Le chiffrement asymétrique opère sur des
entiers plus petits que $n$ : il faudrait découper le message en blocs, et chaque
bloc identique produirait le même chiffré, ce qui laisserait apparaître la
structure du message. Les schémas symétriques sont conçus pour traiter des flux
de longueur quelconque sans cette faiblesse.

C'est pourquoi HTTPS combine les deux : l'asymétrique une fois, pour convenir
d'un secret ; le symétrique ensuite, pour tout le reste de la conversation.
\end{corrige}
